\subsection{三维贯通实验：完整参数表与逐层中间值}
\phantomsection\label{app:transformer-intermediates}

正文\hyperref[intro:22b-token-to-output]{三维贯通章节}只保留形状和代表元素；
这里保存同一次前向传播的所有矩阵与完整数值，
方便读者逐项验算。数值来自
\texttt{scripts/teaching\_reference.py} 的 \texttt{parameters()}，
显示时保留四位小数，运算应使用未舍入数。

\textbf{输入与位置。}词表大小 $5$，隐藏宽度 $3$，
输入 ID $[1,3,4]$，目标 ID $[3,4,2]$：
\[
E=\begin{bmatrix}
0&0&0\\1.2&0&-.2\\.2&-.4&.7\\-.3&1.1&.2\\.15&-.2&1.25
\end{bmatrix}\;(5\times3),\qquad
P=\begin{bmatrix}
.10&.05&-.10\\0&-.05&.12\\-.12&.10&.07
\end{bmatrix}\;(3\times3).
\]
第一步为 $X=E_{[1,3,4]}+P$，逐行 LayerNorm
采用 $\gamma=1$、$\beta=0$、$\varepsilon=10^{-5}$：
\[
U=\operatorname{LN}(X)\approx
\begin{bmatrix}
1.3833&-.4368&-.9464\\
-1.1901&1.2566&-.0665\\
-.6033&-.8061&1.4093
\end{bmatrix}.
\]

\textbf{互不相同的 Query、Key 与 Value 投影。}
\[
\begin{aligned}
W_Q&=\begin{bmatrix}.8&.2&-.1\\.1&.9&.25\\-.1&.15&.8\end{bmatrix},\\[3pt]
W_K&=\begin{bmatrix}.85&0&.2\\-.1&.9&0\\.15&.2&.75\end{bmatrix},\\[3pt]
W_V&=\begin{bmatrix}1&.2&0\\0&.9&.15\\.1&-.2&.85\end{bmatrix}.
\end{aligned}
\]
三个投影都为 $3\times3$，对应的输出是
\[
\begin{aligned}
Q&=UW_Q\approx
\begin{bmatrix}
1.1576&-.2584&-1.0047\\
-.8198&.8829&.3800\\
-.7041&-.6347&.9863
\end{bmatrix},\\[4pt]
K&=UW_K\approx
\begin{bmatrix}
1.0775&-.5824&-.4332\\
-1.1473&1.1176&-.2879\\
-.2208&-.4436&.9364
\end{bmatrix},\\[4pt]
V&=UW_V\approx
\begin{bmatrix}
1.2886&.0728&-.8700\\
-1.1968&.9062&.1320\\
-.4623&-1.1280&1.0770
\end{bmatrix}.
\end{aligned}
\]

\textbf{因果分数与残差。}先除以 $\sqrt3$，
将未来位置掩为 $-\infty$：
\[
S=QK^\top/\sqrt3+C\approx
\begin{bmatrix}
1.0583&-\infty&-\infty\\
-.9019&1.0496&-\infty\\
-.4713&-.1071&.7855
\end{bmatrix}.
\]
正文的 $A=\operatorname{softmax}(S)$
保留四位小数；随后
\[
R_{\rm raw}=AV\approx
\begin{bmatrix}
1.2886&.0728&-.8700\\
-.8876&.8025&.0074\\
-.3458&-.4345&.5215
\end{bmatrix},
\quad
W_O=\begin{bmatrix}
.9&0&.1\\.1&.85&0\\0&.15&.95
\end{bmatrix},
\]
\[
X_{\rm att}=X+R_{\rm raw}W_O\approx
\begin{bmatrix}
2.4670&-.0186&-.9976\\
-1.0186&1.7333&.2382\\
-.3247&-.3911&1.7809
\end{bmatrix}.
\]

\textbf{第二次归一化与 FFN。}网络隐藏层宽度为 $4$：
\[
W_{\rm up}=\begin{bmatrix}
.4&.2&-.3&.1\\
-.1&.5&.15&-.25\\
.2&.1&.45&.3
\end{bmatrix}\;(3\times4),
\quad
W_{\rm down}=\begin{bmatrix}
.25&-.1&.15\\
-.05&.2&.1\\
.1&.15&.15\\
.1&-.12&.18
\end{bmatrix}\;(4\times3).
\]
\[
B=\operatorname{LN}(X_{\rm att})\approx
\begin{bmatrix}
1.3601&-.3444&-1.0157\\
-1.1879&1.2585&-.0706\\
-.6739&-.7398&1.4137
\end{bmatrix},
\]
\[
F_{\rm mid}=\operatorname{ReLU}(BW_{\rm up})\approx
\begin{bmatrix}
.3753&0&0&0\\
0&.3846&.5134&0\\
.0872&0&.7274&.5417
\end{bmatrix}.
\]
$H=X_{\rm att}+F_{\rm mid}W_{\rm down}$ 为 $3\times3$，完整数值见本节末。

\textbf{输出头与 logits。}
\[
W_{\rm out}=\begin{bmatrix}
0&.2&0&.2&-.2\\
0&.2&0&1.2&-.4\\
0&-.1&.4&.1&.2
\end{bmatrix}\;(3\times5),
\]
\[
Z=HW_{\rm out}\approx
\begin{bmatrix}
0&.5951&-.3765&.3507&-.6780\\
0&.1448&.1415&2.1027&-.4868\\
0&-.3064&.8002&-.2620&.5776
\end{bmatrix}.
\]
正文的 $\widehat P=\operatorname{softmax}(Z)$ 为 $3\times5$，
目标 $[3,4,2]$ 的三行概率平均交叉熵约 $1.798430$。


\subsubsection*{完整前向矩阵（第 23 章移入）}
下列矩阵均来自同一次前向计算；显示时保留四位小数，计算请使用未舍入参数。
\noindent\textbf{输入 $X$。}\par
\begin{equation}
\boxed{\begin{aligned}
X_{1,:}&=(1.30,\ 0.05,\ -0.30),\\
X_{2,:}&=(-0.30,\ 1.05,\ 0.32),\\
X_{3,:}&=(0.03,\ -0.10,\ 1.32).
\end{aligned}}
\quad X\in\R^{3\times3}.

\end{equation}

\noindent\textbf{注意力 $A$。}\par
\begin{equation}
\boxed{A\approx\begin{bmatrix}
1&0&0\\
0.1244&0.8756&0\\
0.1680&0.2418&0.5903
\end{bmatrix}.}

\end{equation}

\noindent\textbf{隐藏输出 $H$。}\par
\begin{equation}
\boxed{H\approx\begin{bmatrix}
2.5609&-0.0561&-0.9413\\
-0.9865&1.8872&0.3537\\
-0.1760&-0.3557&2.0005
\end{bmatrix}\in\R^{3\times3}.}

\end{equation}

\noindent\textbf{词表概率 $\widehat P$。}\par
\begin{equation}
\boxed{\widehat P\approx\begin{bmatrix}
0.1843&0.3341&0.1264&0.2617&0.0935\\
0.0826&0.0954&0.0951&0.6761&0.0507\\
0.1535&0.1130&0.3418&0.1181&0.2735
\end{bmatrix}\in\R^{3\times5}.}

\end{equation}

\subsubsection*{输出头梯度矩阵（第 24 章移入）}
正文只保留代表梯度与形状，这里保留十五个元素供逐项复算。
\[
\frac{\partial L}{\partial W_{\rm out}}
\approx\begin{bmatrix}
0.1211&0.2472&0.1153&-0.8595&0.3760\\
0.0303&0.0404&0.1355&0.4251&-0.6313\\
0.0543&-0.0182&-0.4674&0.3902&0.0411
\end{bmatrix}.
\]


\begin{idea}
矩阵形状检查顺序：
$E_{\mathrm{ids}}(3\times3)\to X(3\times3)\to
Q,K,V(3\times3)\to A(3\times3)\to H(3\times3)
\to Z(3\times5)\to\widehat P(3\times5)$。
正文负责解释\textbf{为什么这样算}，本附录负责回答\textbf{数字从哪里来}。
\end{idea}
